\scanpage{97}
\mspage{12}
% archive sheet number: 96

3. \uline{Cas} \deleted{où} \inserted{$X$ \uline{régulier et}} $Y=X-U$ \deleted{\unsure{est}} \uline{un diviseur à croisements normaux.} \inserted{\uline{Enoncé du th. principal et de ses corollaires}}

3.1 \inserted{On suppose $\ell\ne p$.}
\deleted{(Dans le \unsure{résultat qui suit}, \unsure{nous n'indiquons que|nous indiquons que} l'action du gr d'inertie, nous supposons \inserted{\unsure{(d. la réd. définit.)}} pour simplifier les notations que $S$ est \uline{local}.)}
(\unsure{Nous} \unsure{noterons} \deleted{cependant dans le cas général} que les hypothèses : $X$ régulier, $Y$ diviseur à croisements normaux, sont stables par \deleted{\unsure{passage de} $\unsure{X}\unsure{\tilde X}$} \inserted{changt. de base $\tilde S\mathbin{-}S$}).

\inserted{Nous supposons de plus que tout schéma régulier \unsure{quasifini} sur $X$ satisfait au th. de pureté cohomologique (cf 4.1).}
% slanted insertion from the left margin
Nous supposons \unsure{pour l'ins. seulement|pour l'instant seulement} que $H^i(U_{\bar K},\mathbb Z/\ell\mathbb Z)$ est fini, ce qui est le cas en tous cas si $U$ est \unsure{propre lisse} sur $K$ (\unsure{est|et} un schéma régulier).
\deleted{\unsure{résultat} \illegible{3} \unsure{de schéma} \illegible{3} $X$.}
\deleted{Nous supposons de plus que \illegible{5} \unsure{l'hypothèse} \illegible{2} $\Lambda_\nu$.}
\deleted{\unsure{rég.} \illegible{3}) \unsure{satisfait au th. de pureté} \illegible{1}}
\margin{\deleted{\unsure{Réd. définitive}}}
\margin{\unsure{de 1 et 3}}

\uline{Théorème} 3.2. Avec les hypothèses du titre, \deleted{et} \deleted{et, pour $i$, $H^i(U_{\bar K},\mathbb Z/\ell\mathbb Z)$ fini} \inserted{de 3.1} \deleted{\unsure{$S$ étant clos séparément}} $H^{\unsure{i|*}}(U_{\bar K},\mathbb Z_\ell)$ est un $\mathbb Z_\ell$-module de type fini et l'action du \inserted{groupe d'inertie} $I$ sur \inserted{$E=H^i(U_{\bar K},\mathbb Z_\ell)^P$} est essentiellement unipotente. \deleted{\unsure{Pour}} ($P$ étant le $p$-sous-groupe de Sylow de $I$, cf I\unsure{\ldots}).

Plus précisément, \deleted{\unsure{soit}} posons (
\[
\deleted{\unsure{X_0}}\qquad X_0=\sum m_{\unsure{\lambda|A}}X_0^{(\lambda)}
\]
\deleted{\unsure{où}}
(identité de diviseurs sur $X$), \deleted{et} où les $X_0^{(\lambda)}$ sont les composantes irréductibles de $X_0$, et les $m_\lambda>0$. Soit $m$ le ppcm des $m_\lambda$, \deleted{et} $m'$ la composante première à $p$ de $m$, et $I'$ le sous-groupe \inserted{ouvert} d'indice $m'$ de $I$. Alors l'action de $I'$ sur $E=\deleted{\unsure{H^i(U_{\bar K},\mathbb Z_\ell)}}\inserted{H^i(U_{\bar K},\mathbb Z_\ell)^P}$ est \uline{unipotente d'échelon $i+1$}, i.e. il y a une filtration \inserted{pure} de $E$ \inserted{par des $\mathrm{Fil}^p$}, invariante sous $I$, ayant $\unsure{p|i}+1$ quotients, et l'action de $I$ sur les $\unsure{\mathrm{Gr}^p|\mathrm{Gr}_F^p}(E)$ étant triviale. En d'autres termes,
